% !TEX root = ../main.tex
\chapter{Kiểm thử và đánh giá}

\section{Mục tiêu và nguyên tắc đánh giá}

Việc đánh giá tập trung vào các rủi ro chính của đề tài:

\begin{itemize}
    \item CV phải giữ cấu trúc A4;
    \item PDF tải lên phải được kiểm tra;
    \item Job không khả dụng không được nhận hồ sơ;
    \item dữ liệu lịch sử không mất khi Job bị xóa;
    \item nội dung rich text và dữ liệu AI không được mang mã thực thi;
    \item kết quả AI phải tuân theo hợp đồng cấu trúc;
    \item trạng thái đọc phải tồn tại qua lần tải lại.
\end{itemize}

Báo cáo phân biệt ba mức bằng chứng:

\begin{itemize}
    \item \textbf{Đã chạy:} kết quả được lấy từ lệnh thực thi trong môi trường đánh giá.
    \item \textbf{Đối chiếu mã nguồn:} xác nhận thành phần hoặc điều kiện tồn tại nhưng chưa chứng minh luồng chạy toàn hệ thống.
    \item \textbf{Kịch bản đề xuất:} kiểm thử thủ công cần thực hiện trước nghiệm thu; báo cáo không tự gán kết quả đạt cho các kịch bản chưa chạy.
\end{itemize}

\section{Môi trường và khả năng tái lập}

\begin{table}[H]
    \centering
    \small
    \begin{tabularx}{\textwidth}{|p{0.27\textwidth}|Y|} \hline
        \textbf{Thông tin} & \textbf{Giá trị ghi nhận} \\ \hline
        Thời điểm đánh giá & Ngày 09/08/2026, múi giờ Việt Nam. \\ \hline
        Mã nguồn sản phẩm & Nhánh \texttt{main}, working tree tại \texttt{D:\textbackslash Github\textbackslash DA2}; commit gần nhất \texttt{dcc8c58}. \\ \hline
        Trạng thái nguồn & Working tree có cả thay đổi đã theo dõi và tệp chưa theo dõi tại thời điểm ghi nhận; riêng hash commit không đủ tái tạo phiên bản đã kiểm tra. \\ \hline
        Runtime & Node.js 26.5.0; npm 11.17.0. \\ \hline
        Chuỗi kiểm tra & \texttt{npm run check}, lần lượt chạy ESLint, Node Test Runner và Vite production build. \\ \hline
        Phạm vi không dùng & Không gọi API AI trả phí để chấm ngữ nghĩa; không chạy tải, pentest hoặc browser end-to-end trong lần đánh giá này. \\ \hline
    \end{tabularx}
    \caption{Môi trường đánh giá mã nguồn DA2}
    \label{tab:test-environment}
\end{table}

Việc ghi rõ working tree chưa sạch là cần thiết: thay đổi hiện tại trải rộng từ dependency, cấu hình Vite/Vercel, giao diện, route, service, model đến test và tài liệu, nhưng chưa thuộc commit \texttt{dcc8c58}. Để tái lập tuyệt đối, phiên bản sản phẩm cần được commit hoặc phát hành bằng tag sau khi loại tệp tạm và bí mật.

\section{Kết quả quy trình kiểm tra tự động}

\begin{table}[H]
    \centering
    \small
    \begin{tabularx}{\textwidth}{|p{0.22\textwidth}|p{0.23\textwidth}|Y|p{0.13\textwidth}|} \hline
        \textbf{Bước} & \textbf{Lệnh con} & \textbf{Kết quả ghi nhận} & \textbf{Đánh giá} \\ \hline
        Phân tích tĩnh & \texttt{npm run lint} & ESLint hoàn tất, 0 lỗi và 0 cảnh báo. & Đạt \\ \hline
        Kiểm thử logic & \texttt{npm test} & 41/41 trường hợp đạt, 0 thất bại. & Đạt \\ \hline
        Đóng gói frontend & \texttt{npm run build} & Vite 7.3.6 build thành công, 1.906 module được biến đổi; không phát cảnh báo kích thước chunk. & Đạt \\ \hline
        Toàn chuỗi & \texttt{npm run check} & Cả ba bước hoàn tất với mã thoát thành công. & Đạt \\ \hline
    \end{tabularx}
    \caption{Kết quả chạy quy trình kiểm tra ngày 09/08/2026}
    \label{tab:check-results}
\end{table}

So với mốc trước, hai cảnh báo React Hook đã được loại bỏ và trang báo cáo quản trị không còn dùng dữ liệu minh họa. Các trang được nạp lười; Vite tách riêng nhóm PDF, rich text, React, giao diện, form, query, router, thời gian thực và ngày giờ. Ba nhóm JavaScript lớn nhất trong output lần này là PDF khoảng 433~kB, rich text khoảng 428~kB và React khoảng 385~kB; PDF worker là tài nguyên riêng khoảng 1,26~MB. Build sạch xác nhận khả năng đóng gói, không tự chứng minh thời gian tải hoặc hiệu năng mạng đạt yêu cầu.

\section{Kiểm kê bộ kiểm thử}

Mười hai tệp test sử dụng Node Test Runner, tổng cộng 41 trường hợp. Nội dung dưới đây được tổng hợp theo đúng tên test trong mã nguồn.

\begin{table}[H]
    \centering
    \footnotesize
    \begin{tabularx}{\textwidth}{|p{0.30\textwidth}|p{0.08\textwidth}|Y|} \hline
        \textbf{Tệp test} & \textbf{Số} & \textbf{Phạm vi chính} \\ \hline
        \texttt{a4Layout.test.js} & 3 & Kích thước A4, hành động tải PDF dựa trên print và render PDF upload không dùng native iframe. \\ \hline
        \texttt{aiDataSanitizer.test.js} & 5 & Loại HTML thực thi, giới hạn nội dung, tách lời khuyên khỏi CV, bỏ thông tin liên hệ và kỹ năng chứa mã nguy hiểm. \\ \hline
        \texttt{aiIntegration.test.js} & 4 & Route AI có xác thực, structured output, các workflow frontend, trang phỏng vấn và điều hướng theo role. \\ \hline
        \texttt{companyUi.test.js} & 3 & Nhãn tiếng Việt, chuẩn hóa thông tin Job và định dạng ngày trên hồ sơ công khai. \\ \hline
        \texttt{dateInput.test.js} & 5 & Chuẩn hóa ngày cũ, ngày Resume, hoạt động đang diễn ra và chuyển đổi biểu diễn tháng. \\ \hline
        \texttt{interviewData}\allowbreak\texttt{Sanitizer.test.js} & 3 & Giới hạn thiết lập, làm sạch lịch sử 12 câu và giới hạn feedback đầu ra. \\ \hline
    \end{tabularx}
    \caption{Danh mục tệp kiểm thử tự động (phần 1)}
    \label{tab:test-inventory-part1}
\end{table}

\begin{table}[H]
    \centering
    \footnotesize
    \begin{tabularx}{\textwidth}{|p{0.30\textwidth}|p{0.08\textwidth}|Y|} \hline
        \textbf{Tệp test} & \textbf{Số} & \textbf{Phạm vi chính} \\ \hline
        \texttt{jobAvailability.test.js} & 3 & Job đã xóa, Job active còn hạn và Job hết hạn. \\ \hline
        \texttt{jobDeadlineUi.test.js} & 3 & Ưu tiên trạng thái hết hạn, chặn ứng tuyển trước khi đồng bộ và cho phép Job còn hạn. \\ \hline
        \texttt{pdfValidation.test.js} & 3 & Nhận PDF data URL, từ chối Word/nội dung giả và chặn file lớn hơn 5~MB. \\ \hline
        \texttt{readStatus}\allowbreak\texttt{Persistence.test.js} & 2 & Lưu trạng thái đọc message và notification ở server thay vì chỉ ẩn badge trong React state. \\ \hline
        \texttt{resumeAiFill.test.js} & 4 & Merge kỹ năng an toàn, chỉ điền mục trống và giữ mô tả kỹ thuật hợp lệ. \\ \hline
        \texttt{richTextSanitizer.test.js} & 3 & Giữ định dạng CV được hỗ trợ, loại HTML nhúng/thực thi và xử lý input không phải chuỗi. \\ \hline
        \textbf{Tổng} & \textbf{41} & \textbf{Toàn bộ 41 trường hợp đạt.} \\ \hline
    \end{tabularx}
    \caption{Danh mục tệp kiểm thử tự động (phần 2)}
    \label{tab:test-inventory}
\end{table}

\section{Truy vết yêu cầu và bằng chứng}

\begin{table}[H]
    \centering
    \footnotesize
    \begin{tabularx}{\textwidth}{|p{0.19\textwidth}|p{0.38\textwidth}|Y|} \hline
        \textbf{Yêu cầu} & \textbf{Bằng chứng tự động} & \textbf{Mức bao phủ} \\ \hline
        FR-CV-01, FR-CV-02 & \texttt{a4Layout}, \texttt{dateInput}, \texttt{richTextSanitizer} & Bao phủ layout, ngày và làm sạch nội dung CV; chưa bao phủ CRUD qua database thật. \\ \hline
        FR-CV-03 & \texttt{pdfValidation}, \texttt{a4Layout} & Bao phủ validator và cơ chế render; chưa đo PDF đa trang trên nhiều trình duyệt. \\ \hline
        FR-AI-01--03 & \texttt{aiDataSanitizer}, \texttt{aiIntegration}, \texttt{resumeAiFill}, \texttt{interviewDataSanitizer} & Bao phủ hợp đồng, sanitizer và wiring; chưa chấm chất lượng ngữ nghĩa từ provider. \\ \hline
        FR-JOB-01, FR-APP-01 & \texttt{jobAvailability}, \texttt{jobDeadlineUi} & Bao phủ khả dụng Job; unique index và transaction cần integration test với MongoDB. \\ \hline
        FR-APP-02 & Đối chiếu model Application và logic \texttt{jobSnapshot}. & Có bằng chứng mã nguồn, chưa có test tự động chuyên biệt. \\ \hline
        FR-COM-01 & \texttt{companyUi} & Bao phủ một phần UI; chưa bao phủ toàn bộ vòng đời Job/Application. \\ \hline
        FR-MSG-01 & \texttt{readStatusPersistence} & Bao phủ trạng thái đọc ở mức mã; chưa kiểm thử quyền phòng, ảnh chat hoặc Socket.IO end-to-end. \\ \hline
        FR-ADM-01 & Điều hướng vai trò trong \texttt{aiIntegration}; ESLint chạy sạch. & Bao phủ hạn chế, cần kiểm thử ma trận quyền và API bằng tài khoản thật. \\ \hline
    \end{tabularx}
    \caption{Ma trận truy vết yêu cầu--kiểm thử}
    \label{tab:requirements-tests}
\end{table}

\clearpage
\section{Kịch bản kiểm thử thủ công cần thực hiện}

Các kịch bản dưới đây là checklist nghiệm thu, chưa được ghi nhận là đã chạy trong phiên đánh giá ngày 09/08/2026.

\begin{table}[H]
    \centering
    \footnotesize
    \begin{tabularx}{\textwidth}{|p{0.20\textwidth}|Y|p{0.19\textwidth}|} \hline
        \textbf{Nhóm} & \textbf{Kịch bản và tiêu chí chấp nhận} & \textbf{Trạng thái} \\ \hline
        Xác thực/RBAC & Đăng nhập bằng ứng viên, doanh nghiệp, admin; thử URL và API sai vai trò; tải lại phiên; đăng xuất. Kiểm tra riêng AI, lưu việc, ứng tuyển và phòng chat ở backend. & Chưa chạy thủ công \\ \hline
        CV/A4/PDF & Tạo CV ngắn và dài, in bằng các trình duyệt mục tiêu, đo trang 210~x~297~mm; mở PDF nhiều trang; thử file giả, sai MIME và quá 5~MB. & Chưa chạy thủ công \\ \hline
        Ứng tuyển & Nộp bằng Resume và PDF, gửi lặp, thử Job draft/expired/deleted; xóa Job rồi kiểm tra lịch sử từ snapshot. & Chưa chạy thủ công \\ \hline
        Doanh nghiệp & Tạo/sửa/đóng/mở/xóa Job, xem CV đúng và sai quan hệ, chuyển toàn bộ trạng thái Application. & Chưa chạy thủ công \\ \hline
        Chat/thông báo & Gửi giữa hai trình duyệt, thử phòng không thuộc người gọi, ảnh lớn/sai loại, mất và kết nối lại mạng, đánh dấu đọc, ẩn phòng một phía. & Chưa chạy thủ công \\ \hline
        AI & Thử năm workflow, lỗi thiếu khóa, timeout, hết quota và payload nguy hiểm; xác nhận AI không tự ghi đè Resume. & Chưa chạy thủ công \\ \hline
    \end{tabularx}
    \caption{Checklist nghiệm thu thủ công chưa thực hiện}
    \label{tab:manual-checklist}
\end{table}

\section{Đánh giá riêng cho chức năng AI}

Kết quả 16 test thuộc bốn tệp tập trung vào AI xác nhận được các thuộc tính có tính quyết định: dữ liệu được làm sạch và giới hạn, route yêu cầu xác thực, nhà cung cấp được yêu cầu trả structured output, kết quả fill không ghi thêm sự kiện vào mục đã có và feedback phỏng vấn bị giới hạn. Ba test rich text bổ sung bằng chứng cho nội dung CV nhưng được tính ở nhóm riêng. Đây là bằng chứng về đường nối và hợp đồng phần mềm, không phải bằng chứng rằng nội dung do mô hình sinh luôn đúng hoặc hữu ích.

Để đánh giá chất lượng ngữ nghĩa, quy trình cần thực hiện các bước sau:

\begin{enumerate}
    \item Xây dựng bộ dữ liệu cố định gồm CV ở nhiều mức kinh nghiệm và mô tả công việc tương ứng.
    \item Ẩn danh thông tin cá nhân trong dữ liệu đánh giá.
    \item Chạy cùng mô hình, prompt và tham số cho toàn bộ mẫu.
    \item Có ít nhất hai người chấm độc lập theo các tiêu chí:
    \begin{itemize}
        \item bám sát dữ liệu nguồn;
        \item tính đúng đắn;
        \item khả năng hành động;
        \item độ rõ ràng;
        \item không bịa đặt;
        \item không tạo phân biệt đối xử.
    \end{itemize}
\end{enumerate}

Phụ lục A cung cấp rubric đề xuất. Chỉ sau bước này mới nên báo cáo điểm trung bình, độ đồng thuận giữa những người chấm hoặc kết quả so sánh mô hình.

\section{Hạn chế và rủi ro còn lại}

\begin{itemize}
    \item Bộ test thiên về unit và source-level integration; chưa chạy MongoDB thật, browser end-to-end, ma trận trình duyệt A4 hoặc Socket.IO khi mạng chập chờn.
    \item Working tree DA2 chưa có commit/tag tái lập; chưa có CI, số liệu bao phủ mã hoặc artifact kiểm thử được lưu theo phiên bản.
    \item Kiểm tra quyền theo vai trò chưa đồng đều: API AI, ứng tuyển và lưu việc chủ yếu dựa trên JWT; gửi/ẩn chat chưa xác minh đầy đủ người gọi thuộc phòng.
    \item Khóa JWT dự phòng và tài khoản quản trị seed chỉ phù hợp phát triển; ảnh chat chưa có validator server tương đương PDF.
    \item Chất lượng ngữ nghĩa, tính công bằng và chi phí AI chưa được chấm trên bộ dữ liệu độc lập; chưa gọi provider thật trong chuỗi kiểm tra tự động.
    \item Chưa thực hiện kiểm thử tải, kiểm thử xâm nhập, CSRF chuyên sâu, accessibility audit, Vercel Cron/Socket.IO production hoặc khôi phục cơ sở dữ liệu và tệp.
\end{itemize}

\noindent\textbf{Kết luận đánh giá.}
Trong phạm vi đã chạy, RabbitCV đạt 41/41 test, ESLint không có lỗi hoặc cảnh báo và build frontend thành công mà không phát cảnh báo kích thước chunk. Đây là baseline kỹ thuật của working tree ngày 09/08/2026, chưa phải bằng chứng sẵn sàng production; kiểm thử thủ công, tích hợp MongoDB, ma trận quyền, hạ tầng triển khai và đánh giá ngữ nghĩa AI vẫn cần hoàn tất.
